
\documentclass[../../../physics-standard_answers_all.tex]{subfiles}

\begin{document}

\setlength{\abovedisplayskip}{2mm}
\setlength{\belowdisplayskip}{1mm}

\begin{enumerate}

    \item まず，$z_1 = \uwave{0}$ であり，キルヒホッフ則より，
        %
        \begin{align*}
            \frac{x_1}{\varepsilon_0 S} \cdot 2d = V_0
            \qquad \text{∴} \; x_1
            = \uwave{\frac{\varepsilon_0 S}{2d} V_0} \;.
        \end{align*}
        %
        すると，
        %
        \begin{align*}
            E_1 = \frac{x_1}{\varepsilon_0 S}
            = \uwave{\frac{V_0}{2d}} \;.
        \end{align*}
        
    \item XとZからなる部分についての電荷保存則より，
        %
        \begin{align*}
            x_2 + z_2 = x_1 + z_1 \;.
            \qquad \text{∴} \; x_2 + z_2 =
            \frac{\varepsilon_0 S}{2d} V_0 \;.
        \end{align*}
        %
        キルヒホッフ則より，
        %
        \begin{align*}
            & \frac{x_2}{\varepsilon_0 S} \cdot 2d
            = \frac{z_2}{\varepsilon_0 S} \cdot d \\
            & \text{∴} \; x_2
            =\uwave{\frac{\varepsilon_0 S}{6d} V_0} \;,
            \quad z_2 = \uwave{\frac{\varepsilon_0 S}{3d} V_0} \;.
        \end{align*}
        %
        すると，
        %
        \begin{align*}
            E_2 = \frac{x_2}{\varepsilon_0 S}
            = \uwave{\frac{V_0}{6d}} \;,  
            \quad \phi_2 = E_2 \cdot 2d = \uwave{\frac{1}{3} V_0} \;.
        \end{align*}

    \item 電荷保存則より，Zの電荷は変わらない．つまり，
        %
        \begin{align*}
            z_3 = z_2 = \uwave{\frac{\varepsilon_0 S}{3d} V_0} \;.
        \end{align*}
        %
        キルヒホッフ則より，
        %
        \begin{align*}
            \frac{x_3}{\varepsilon_0 S} \cdot 2d = V_0
            \qquad \text{∴} \; x_3
            = \uwave{\frac{\varepsilon_0 S}{2d} V_0} \;.
        \end{align*}

    \item キルヒホッフの法則より，
        %
        \begin{align*}
            & \frac{x_4}{\varepsilon_0 S} \cdot 2d
            = \frac{z_4}{\varepsilon_0 S} \cdot d
            = V_0 \;. \\ 
            & \text{∴} \; x_4
            = \uwave{\frac{\varepsilon_0 S}{2d} V_0} \;,
            \quad z_4 = \uwave{\frac{\varepsilon_0 S}{d} V_0} \;.
        \end{align*}

    \item XとZからなる部分についての電荷保存則より，
        %
        \begin{align*}
            x_5 + z_5 = x_4 + z_4 \;.
            \qquad \text{∴} \; x_5 + z_5 =
            \frac{3\varepsilon_0 S}{2d} V_0 \;.
        \end{align*}
        %
        キルヒホッフ則より，
        %
        \begin{align*}
             \frac{x_5}{\varepsilon_0 S} \cdot 2d
            = \frac{z_5}{\varepsilon_0 S} \cdot 2d
            \qquad \text{∴} \; x_5 = z_5
            =\uwave{\frac{3\varepsilon_0 S}{4d} V_0} \;.
        \end{align*}

\end{enumerate}

\end{document}
